% =============================================================
%  00b_ambiente.tex — Guida rapida ad Allegro CL Free Express
%  Concisa e pratica: avvio, REPL, comandi essenziali
% =============================================================
\chapter*{Usare Allegro CL Free Express}
\addcontentsline{toc}{chapter}{Usare Allegro CL Free Express}
\markboth{Usare Allegro CL Free Express}{Usare Allegro CL Free Express}

% --- Colori per questa sezione (riusa quelli del cheatsheet) ---
\colorlet{envBlue}{csBlue}
\colorlet{envGreen}{csGreen}
\colorlet{envOrange}{csOrange}
\colorlet{envPurple}{csPurple}

% --- Box per le sezioni ambiente ---
% Unbreakable, testo grigio/piccolo, codice in evidenza
\newtcolorbox{envbox}[2][]{%
  enhanced,
  colback=#2!4,
  colframe=#2!70!black,
  fonttitle=\bfseries\sffamily,
  coltitle=white,
  title={#1},
  boxrule=0.8pt,
  arc=4pt,
  left=12pt, right=12pt, top=10pt, bottom=10pt,
  toptitle=4pt, bottomtitle=4pt,
  before skip=14pt plus 4pt,
  after skip=14pt plus 4pt,
  fontupper=\small\color{black!75},
  parskip=6pt
}

% --- Box per suggerimenti ---
\newtcolorbox{tipbox}[1][]{%
  enhanced,
  colback=envPurple!6,
  colframe=envPurple!50!black,
  fonttitle=\bfseries\sffamily\small,
  coltitle=white,
  title={#1},
  boxrule=0.5pt,
  arc=3pt,
  left=10pt, right=10pt, top=6pt, bottom=6pt,
  before skip=10pt,
  after skip=10pt,
  fontupper=\small\color{black!70}
}

% --- Box per attenzione/warning ---
\newtcolorbox{warnbox}[1][]{%
  enhanced,
  colback=envOrange!6,
  colframe=envOrange!70!black,
  fonttitle=\bfseries\sffamily\small,
  coltitle=white,
  title={#1},
  boxrule=0.5pt,
  arc=3pt,
  left=10pt, right=10pt, top=6pt, bottom=6pt,
  before skip=10pt,
  after skip=10pt,
  fontupper=\small\color{black!70}
}

\vspace{-0.3cm}

\noindent\small\color{black!70}
Allegro CL Free Express (Franz Inc.) è un'implementazione ANSI Common Lisp
gratuita per uso didattico. Tutto il codice di questo manuale è testato con questa edizione.
\color{black}

\vspace{10pt}

% =============================================================
%  AVVIO E REPL
% =============================================================

\begin{envbox}[Avvio e REPL]{envBlue}

Avvia \fn{allegro-ansi.exe} (modalità ANSI: i simboli diventano maiuscoli).
Compare il prompt REPL:\nopagebreak

\begin{lstlisting}
CL-USER(1):
\end{lstlisting}

Digita un'espressione e premi Invio:\nopagebreak

\begin{lstlisting}
CL-USER(1): (+ 1 2)
3
CL-USER(2): (defun quadrato (x) (* x x))
QUADRATO
CL-USER(3): (quadrato 5)
25
\end{lstlisting}

\textbf{Riutilizzare risultati:} \fn{*}, \fn{**}, \fn{***} contengono
gli ultimi tre valori:\nopagebreak

\begin{lstlisting}
CL-USER(4): (+ 10 20)
30
CL-USER(5): (1+ *)       ; * vale 30 -> 31
31
\end{lstlisting}

\end{envbox}

% =============================================================
%  CARICARE E COMPILARE FILE
% =============================================================

\begin{envbox}[Caricare e compilare file]{envGreen}

Scrivi il codice in un file \fn{.lisp}, poi caricalo nel REPL:

\smallskip
\renewcommand{\arraystretch}{1.3}
\begin{tabular}{@{}p{5cm}p{8cm}@{}}
\rowcolor{envGreen!10}
\textbf{Comando} & \textbf{Effetto} \\
\fn{:ld "file.lisp"} & Carica e valuta (interpretato) \\
\rowcolor{envGreen!5}
\fn{:cl "file.lisp"} & Compila e carica (genera \fn{.fasl}) \\
\fn{:cf "file.lisp"} & Solo compila \\
\end{tabular}
\renewcommand{\arraystretch}{1.0}

\smallskip
Gestione directory:\nopagebreak

\begin{lstlisting}
CL-USER(1): :pwd                   ; directory corrente
CL-USER(2): :cd "c:/esercizi-lisp" ; cambia directory
CL-USER(3): :ld "es01.lisp"        ; carica file
\end{lstlisting}

\begin{warnbox}[Percorsi Windows]
Usa slash \fn{/} nei percorsi, non backslash.
Corretto: \fn{"c:/cartella/file.lisp"}
\end{warnbox}

\end{envbox}

% =============================================================
%  COMANDI TOP-LEVEL
% =============================================================

\begin{envbox}[Comandi top-level]{envBlue}

I comandi \fn{:} sono specifici di Allegro (non ANSI standard).

\smallskip
\renewcommand{\arraystretch}{1.2}
\begin{tabular}{@{}l@{\hspace{6pt}}l@{\hspace{10pt}}l@{}}
\rowcolor{envBlue!10}
\textbf{Comando} & \textbf{Abbr.} & \textbf{Effetto} \\
\fn{:ld "file"} & --- & carica file \\
\rowcolor{envBlue!5}
\fn{:cl "file"} & --- & compila e carica \\
\fn{:pwd} / \fn{:cd "dir"} & \fn{:pw} & directory \\
\rowcolor{envBlue!5}
\fn{:trace fn} & \fn{:tr} & attiva trace \\
\fn{:untrace fn} & \fn{:untr} & disattiva trace \\
\rowcolor{envBlue!5}
\fn{:inspect obj} & \fn{:i} & ispeziona oggetto \\
\fn{:describe obj} & \fn{:de} & documentazione \\
\rowcolor{envBlue!5}
\fn{:apropos "str"} & \fn{:ap} & cerca simboli \\
\fn{:history} & \fn{:his} & cronologia \\
\rowcolor{envBlue!5}
\fn{:help} & \fn{:he} & lista comandi \\
\fn{:exit} & \fn{:ex} & esci \\
\end{tabular}
\renewcommand{\arraystretch}{1.0}

\smallskip
\textbf{Cercare funzioni} con \fn{:apropos}:\nopagebreak

\begin{lstlisting}
CL-USER(1): :apropos "remove"
REMOVE           function
REMOVE-IF        function
REMOVE-DUPLICATES function
\end{lstlisting}

\end{envbox}

% =============================================================
%  TRACE
% =============================================================

\begin{envbox}[Trace: seguire la ricorsione]{envOrange}

\fn{:trace} mostra ogni chiamata e il valore restituito:\nopagebreak

\begin{lstlisting}
CL-USER(1): (defun fatt (n)
              (if (= n 0) 1 (* n (fatt (1- n)))))
FATT
CL-USER(2): :trace fatt
CL-USER(3): (fatt 4)
  0: (FATT 4)
    1: (FATT 3)
      2: (FATT 2)
        3: (FATT 1)
          4: (FATT 0)
          4: FATT returned 1
        3: FATT returned 1
      2: FATT returned 2
    1: FATT returned 6
  0: FATT returned 24
24
\end{lstlisting}

Disattivare: \fn{:untrace fatt} (o \fn{:untrace} per tutte).

\end{envbox}

% =============================================================
%  DEBUGGER
% =============================================================

\begin{envbox}[Il debugger]{envOrange}

In caso di errore, Allegro entra nel debugger (prompt \fn{[1]}):\nopagebreak

\begin{lstlisting}
CL-USER(1): (/ 1 0)
Error: Division-by-zero.
[1] CL-USER(2):
\end{lstlisting}

\smallskip
\renewcommand{\arraystretch}{1.2}
\begin{tabular}{@{}l@{\hspace{6pt}}l@{\hspace{10pt}}l@{}}
\rowcolor{envOrange!10}
\textbf{Comando} & \textbf{Abbr.} & \textbf{Effetto} \\
\fn{:reset} & \fn{:res} & torna al livello 0 \\
\rowcolor{envOrange!5}
\fn{:pop} & --- & sale di un livello \\
\fn{:bt} / \fn{:zoom} & \fn{:zo} & stampa stack \\
\rowcolor{envOrange!5}
\fn{:error} & \fn{:err} & ristampa errore \\
\end{tabular}
\renewcommand{\arraystretch}{1.0}

\begin{warnbox}[Livelli annidati]
Se il prompt mostra \fn{[2]} o più, hai accumulato errori.
Usa \fn{:reset} per tornare al livello 0.
\end{warnbox}

\end{envbox}

% =============================================================
%  ERRORI COMUNI
% =============================================================

\begin{envbox}[Errori comuni]{envOrange}

\textbf{Parentesi mancante:} se Invio non fa nulla, manca una \fn{)}.
Chiudila o premi \textbf{Ctrl+C} per annullare.

\smallskip
\textbf{Parentesi in più:}\nopagebreak
\begin{lstlisting}
CL-USER(1): (+ 1 2))
Error: eof encountered ...
\end{lstlisting}
Usa \fn{:reset} e riscrivi.

\smallskip
\textbf{Funzione non definita:}\nopagebreak
\begin{lstlisting}
CL-USER(1): (mia-funzione 3)
Error: Undefined function: MIA-FUNZIONE
\end{lstlisting}
Verifica di aver caricato il file con \fn{:ld}.

\end{envbox}

% =============================================================
%  FLUSSO DI LAVORO
% =============================================================

\begin{envbox}[Flusso di lavoro consigliato]{envPurple}

\begin{enumerate}[leftmargin=1.5em, itemsep=3pt, topsep=0pt]
  \item Crea \fn{es01.lisp} con un editor
  \item Scrivi \textbf{una sola funzione} per volta
  \item Nel REPL: \fn{:cd} + \fn{:ld "es01.lisp"}
  \item Testa subito con vari input
  \item Bug? Modifica, salva, ricarica con \fn{:ld}
  \item Per ricorsione: usa \fn{:trace}
  \item Passa alla prossima solo quando funziona
\end{enumerate}

\begin{tipbox}[Sviluppo incrementale]
Non scrivere mai tutto insieme. Scrivi poco, testa subito, procedi.
\end{tipbox}

\end{envbox}
